\chapter{Literature review}\label{cap:literature}

\section{The dentate gyrus and pattern separation}

% TODO: insert figure

Two properties are usually thought to underlie pattern separation in the DG\@: the expansion of a relatively small entorhinal
population onto a much larger number of GCs, and the competitive network dynamics that keep only a few of those GCs active at any
one time~\parencite{hainmuellerDentate2020}. The sparseness of GC activity has been confirmed by recordings from identified cells:
fewer than 5\% of a rodent's GCs are active while it explores a given environment, and those that are active typically fire
in a single location~\parencite{diamantakiSparse2016, hainmuellerDentate2020}.

This sparseness is not a property of the GCs alone but of the competitive dynamics of the network, in which active GC ensembles
suppress competing ensembles by recruiting disynaptic inhibition~\parencite{hainmuellerDentate2020}. The DG contains several
GABAergic interneuron types that implement this inhibition~\parencite{degroInterneuron2022}. The best characterized are the
parvalbumin-expressing cells, which include the basket cells and inhibit the perisomatic domain of the GCs, and the
somatostatin-expressing cells, which target the distal GC dendrites near their PP input
synapses~\parencite{hainmuellerDentate2020}. The functional consequence of losing this balance is clear: when a small number
of GCs become abnormally dominant, the discriminative function of the region degrades~\parencite{scharfmanCorruption2016}, and the
spatial reach of the inhibition determines how locally the network synchronizes~\parencite{struberDistancedependent2017}.

Two further features of the DG architecture matter for how this inhibition is organized. The first is the lamellar arrangement of
the region, in which the GC population and its projections are organized into transverse slices that operate with partial
independence~\parencite{sloviterUpdating2012}, although the interlamellar connectome is denser than the strict form of the
hypothesis assumed~\parencite{pakHippocampal2022}. The second is the mossy cell population of the hilus, which receives excitation
from the GCs and projects back to GCs and to hilar interneurons, so that its net effect on the GC population is
inhibitory~\parencite{scharfmanHilar2013}. Mossy cells contribute to pattern separation in their own
right~\parencite{goodsmithDentate2019}, and modeling work attributes that contribution specifically to the indirect feedback
inhibition they recruit~\parencite{galloniRecurrent2022}.

What is less settled is the form the separation takes. Morphing the enclosure through a series of intermediate shapes,
\textcite{leutgebPattern2007} found that the firing fields of a single granule cell changed noncoherently with one another,
decorrelating the DG population strongly even between the closest shapes, while no new granule cells were recruited as the
environments were made more different, while CA3 behaved in the opposite way. Separation in the DG is therefore a redistribution
of firing rates within an already active population rather than a turnover of which cells are active. How completely it separates
is nonetheless qualified by the finding that a substantial fraction of mature GCs encodes features common to several
contexts~\parencite{hainmuellerDentate2020}, and reviews treat pattern separation as one of several operations the region
supports, alongside novelty detection and the binding of information to spatial context~\parencite{kesnerAnalysis2018}.

\section{CA3: autoassociation and pattern completion}

The recurrent collaterals of CA3 give the region attractor dynamics: an input resembling a stored pattern drives the network
toward the stored state, which is the mechanism usually invoked to explain pattern completion~\parencite{rollsMechanisms2013,
rollsStorage2018, miglioreCA32019}. The same dynamics can also push CA3 into a different attractor when its inputs are
sufficiently distinct from one another, so the output of the region is not simply a smoothed version of its
input~\parencite{hainmuellerDentate2020}.

The relative weight of the two inputs to CA3 is worth noting, because it sets a limit on how much influence the DG can exert. An
individual pCA3 cell receives inputs from thousands of EC and other pCA3 cells but from only some fifty GCs, and many of those GCs
are silent during exploration~\parencite{hainmuellerDentate2020}. The influence of the GCs is therefore not simply a matter of
excitatory drive: GC burst activity reliably recruits pCA3 cells, whereas low-frequency GC firing preferentially engages
feedforward interneurons and so exerts a net inhibitory effect on CA3~\parencite{hainmuellerDentate2020}. This asymmetry is the
mechanism by which a small GC population can shape CA3 activity without dominating it.

CA3 is not the end of the circuit as far as the DG is concerned. A portion of the pCA3 population projects back to the hilus, a
backprojection that contributes to pattern separation in the DG itself~\parencite{myersPattern2011}. Recent anatomical work
continues to refine the picture: connectivity between pCA3 cells is organized by cell subtype rather than at
random~\parencite{sammonsSubtype2025}.

The DG and CA3 therefore operate as a single loop rather than as two stages in sequence, so that whatever alters the sparseness of
the GC code is felt in CA3 and returns to the DG through the backprojection. The population that performs the DG computation is
not fixed either, since a fraction of it is replaced throughout life, which is the subject of the next section.

\section{Adult hippocampal neurogenesis}

That new neurons are added to the adult hippocampus was first reported from autoradiographic evidence in the
rat~\parencite{altmanAutoradiographic1965}, although the finding met resistance for decades before becoming established in
rodents~\parencite{mingAdult2011, abbottAdult2020}. At least two niches are well regarded to sustain this process in the rodent
brain: the subventricular zone, whose progeny migrate to the olfactory bulb and differentiate into inhibitory neurons, and the
subgranular zone of the DG, which is the only one that generates excitatory glutamatergic cells and therefore the only one that
adds principal neurons to a cortical circuit~\parencite{folszAdultborn2023, faresNeurogenesis2019}.

The rate of addition is low. Of all DG GCs, on the order of 0.2\% are generated daily in rats and 0.06\% in
mice~\parencite{cameronAdult2001, folszAdultborn2023}, and roughly half of the cells produced are eliminated by waves of
programmed cell death during maturation~\parencite{folszAdultborn2023}. Because this addition is balanced by the continuous
removal of mostly perinatally generated mGCs, the total number of GCs stays approximately constant~\parencite{folszAdultborn2023}.
The cells that do survive are stably maintained and eventually become indistinguishable from developmentally born
GCs~\parencite{denoth-lippunerFormation2021, folszAdultborn2023}.

Maturing cells extend their dendrites outward through the molecular layer and send an axon toward
CA3~\parencite{zhaoDistinct2006}, and their electrophysiological properties change in step, recapitulating aspects of embryonic
development~\parencite{espositoNeuronal2005}. Roughly four to seven weeks after birth the cells pass through a critical window in
which they are already wired into the circuit, receiving afferent input and projecting to CA3, but have not yet acquired the
properties of a mature cell. Within that window they are more excitable than mGCs and also more plastic, showing a lower threshold
for the induction of long-term potentiation and a larger potentiation amplitude~\parencite{aimoneRegulation2014,
folszAdultborn2023}. It is this combination, an integrated cell that still fires more readily than its neighbors, that makes the
window the period of interest for the circuit-level questions treated here. The boundaries of this window, however, are not
settled: birthdating studies combined with immediate-early gene expression report that adult-born cells can remain comparatively
excitable for several months, which challenges the view that their contribution is confined to their first few
weeks~\parencite{folszAdultborn2023}.

\section{Functional roles of adult-born granule cells}

The two views of what adult-born cells do, set out in Chapter~\ref{cap:introduction}, differ in whether the cells are treated as
part of the code or as something acting on it. The encoding view is supported by the observation that hyperexcitable DG cells are
preferentially recruited into engrams, so that cells in the critical window are more readily incorporated into any memory trace
than their mature neighbors, and by reports that they are activated during both encoding and
retrieval~\parencite{piattiNeurogenesis2013, folszAdultborn2023}. The modulatory view rests instead on what these cells do to the
rest of the circuit, and it is the one for which circuit-level evidence has accumulated most
recently~\parencite{aimoneComputational2016, mchughAdultborn2022, folszAdultborn2023}.

Three mechanisms have been proposed for that modulatory action, and they are not exclusive. The first is competition for synaptic
partners. Electron microscopy and live imaging of synaptogenesis on adult-born cells showed that dendritic filopodia are most
abundant early in maturation and that the tip of every filopodium examined lay within \SI{200}{\nano\metre} of a pre-existing bouton
already synapsing on another neuron, with new spines contacting multiple-synapse boutons preferentially; the connectivity of these
cells continued to change for at least two months~\parencite{toniSynapse2007}. New cells therefore insert themselves into
established connections rather than building fresh ones, and later displace mGC axons
entirely~\parencite{toniSynapse2007, folszAdultborn2023}.

The second is inhibition of mGCs, both directly and through interneurons. Here the sign of the interaction depends on which input
stream drives the adult-born cell: those receiving lateral entorhinal input inhibit mGCs, whereas those receiving medial
entorhinal input excite them~\parencite{lunaAdultborn2019}. Evidence that the net effect on the surrounding network is inhibitory
comes from removing the cells: after focal irradiation or pharmacogenetic ablation of the young adult-born population,
perforant-path evoked responses in the DG were reduced, but spontaneous gamma-frequency bursts in the DG and hilus grew markedly
in amplitude and dentate firing became more tightly synchronized to them, which led the authors to conclude that young neurons
normally inhibit or destabilize recurrent activity in the DG and hilus~\parencite{lacefieldEffects2012}.

The third is inhibition reaching beyond the DG\@. Adult-born cells engage hippocampal interneurons that inhibit neighboring mGCs
and also exert inhibition on CA3 and CA1, with the reported consequence being an overall sparsification of population firing
across the hippocampus~\parencite{folszAdultborn2023, mchughAdultborn2022}. The clearest evidence that the effect crosses regions
is that suppressing adult neurogenesis, either chemically or genetically, impairs the population-based coding of similar contexts
in CA3 while leaving the coding of dissimilar contexts intact~\parencite{niiboriSuppression2012}. The deficit therefore appears
downstream of the region in which the new cells are, and only for the inputs that are hardest to tell apart. Accounting for both
features in terms of the circuit is the problem this work addresses.

The behavioral evidence is broadly consistent with a role in discrimination but is not uniform. Increasing neurogenesis is
sufficient to improve behavioral pattern separation~\parencite{sahayIncreasing2011}, and cells in the critical window are
required for tasks that demand the separation of highly similar contexts while appearing dispensable for learning a single
location in the water maze~\parencite{folszAdultborn2023}.

A separate proposal is that adult-born cells provide a time stamp, linking memories encoded close together in time while
separating remote ones, on the grounds that the population available for recruitment changes as cells are added and
mature~\parencite{aimoneComputational2009}. This has some \textit{in vivo} support: recordings from the DG identified cells whose
activity was selective for a single context when the exposures to the two contexts were widely separated in time, with that
selectivity attenuating as the exposures were brought closer together and attenuating further when neurogenesis was reduced, which
the authors interpret as a temporally orthogonalizing code in the DG~\parencite{rangelTemporally2014}. Support for the idea is
nonetheless mixed, since little overlap has been found between the adult-born cells active during encoding and during retrieval of
contextual fear memories, which has been taken to suggest that these cells respond to behavioral state rather than to particular
events~\parencite{folszAdultborn2023}. Beyond encoding, adult-born cells have been implicated in later stages of memory
processing, including consolidation during sleep, and in the gradual degradation of old connections that promotes generalization
and eventually forgetting~\parencite{chavanMemory2025, folszAdultborn2023}.

Consistent with the latter, \textcite{mchughAdultborn2022} showed \textit{in vivo} that immature adult-born GCs carry little
spatial information, especially when compared to mGCs; this study also showed that briefly activating those iGCs recruits
inhibitory interneurons throughout the hippocampus and increases population sparsity in both the DG and CA3, pointing to a
circuit-level mechanism by which a small population of new cells could regulate activity well beyond the DG\@. Similarly,
\textcite{koSystems2025} showed that adult neurogenesis drives the reorganization of hippocampal engram circuitry as memories are
consolidated, remodeling the DG-CA3-CA1 network so that its representations generalize over time. These results show that iGCs
have important effects across the hippocampus network, not being restricted to a local function in the DG only.

\section{Computational models of the hippocampus and of adult neurogenesis}

\subsection{Neuron and synapse models}

Any network model has to choose how much of single-cell biophysics to keep. The integrate-and-fire abstraction, which reduces a
neuron to a leaky capacitor with a threshold, dates to the beginning of the last century~\parencite{lapicqueRecherches1907,
abbottLapicque1999} and remains cheap enough for large networks, but it discards the mechanisms that generate the action
potential. At the other extreme, the conductance-based formulation derived from the squid axon~\parencite{hodgkinQuantitative1952}
reproduces those mechanisms explicitly at a substantially higher computational cost. Between the two sit reduced models with a
small number of parameters that nonetheless recover a wide range of firing patterns, of which the Izhikevich model is the one
adopted here~\parencite{izhikevichSimple2003, izhikevichDynamical2006}. How faithfully such reduced models reproduce the spiking
of real neurons has been assessed systematically, and generalized leaky integrate-and-fire variants can classify multiple neuron
types from their responses~\parencite{gerstnerHow2009, teeterGeneralized2018}. The relevant synaptic phenomena for this work are
the short-term ones, in which transmission is depressed by the depletion of releasable resources and facilitated by residual
calcium, and so the Tsodyks--Markram model was used~\parencite{tsodyksNeural1998, zuckerShortTerm2002}. The particular equations
used are stated in Chapter~\ref{cap:methods}.

\subsection{Network models of the dentate gyrus and CA3}

Models of the DG span the same range of detail, from abstract treatments of its coding properties through phenomenological circuit
models to fully conductance-based reconstructions~\parencite{aimoneComputational2016}. Which features prove necessary depends on
the question: the dendritic morphology of the GCs, for instance, turns out to control DG sparsity and separation directly and
cannot simply be abstracted away~\parencite{chavlisDendrites2017}. Models of CA3 range from detailed and simplified pyramidal cell
reconstructions~\parencite{miglioreCA32019} to full-scale data-driven networks in which cell assemblies form and are
retrieved~\parencite{kopsickFormation2024}; broader surveys of hippocampal spiking models and of the microcircuit data that
constrain them are available~\parencite{cutsuridisHippocampal2019, suttonSpiking2021}, as are neuromorphic implementations of the
circuit~\parencite{casanueva-moratoBioinspired2023}.

Two practical constraints shape this literature. The first is that assembling a biologically constrained network requires a
consolidated source of neuron and synapse parameters, which is the role played by knowledge bases such as
Hippocampome.org~\parencite{wheelerHippocampomeorg2023}. The second is computational: simulating large spiking networks on
general-purpose hardware has known limits~\parencite{zenkeLimits2014}, which is one reason models are built at reduced scale and
with efficient simulators~\parencite{stimbergBrian2019a}. It is also worth noting an anatomical fact that constrains any DG-CA3
model: the mossy fiber forms giant multisynapse terminals on its CA3 targets, yet each fiber contacts only on the order of a dozen
pyramidal cells, making it one of the sparsest connections in the brain, and its strength is offset by the heavy tonic inhibition
the GCs receive~\parencite{aimoneComputational2016}.

\subsection{Models of adult neurogenesis}

Computational work on DG neurogenesis has been organized at three scales, namely the individual new neuron, the DG as a circuit,
and the hippocampus as a loop, and the questions asked differ at each~\parencite{aimoneComputational2016, aimoneModeling2011}. At
the level of the single cell, the excitability of the immature neuron has been attributed to a weaker and even depolarizing
response to GABA in the first weeks, to a distinct complement of ion channels giving a higher input resistance, and to the reduced
synaptic input the cell receives; modeling indicates that the reduced input alone is sufficient to bias an immature GC toward
firing more often than a mature one~\parencite{aimoneComputational2016}. That last point matters for the present work, because
reduced afferent input is the mechanism by which immaturity is represented in the model described in Chapter~\ref{cap:methods}.

At the level of the whole hippocampal loop, the earliest models incorporated new neurons into a four-layer circuit and found that
turnover was necessary for effective coding of representationally similar inputs, while later work treating neurogenesis as an
additive rather than a replacement process found that the new cells came to represent much of the response to novel inputs and, in
that configuration, showed signs of impaired pattern separation~\parencite{aimoneComputational2016}. Incorporating neurogenesis into
the most detailed models is difficult for a simple reason: neurogenesis takes place over weeks, whereas conductance-based
models are integrated at sub-millisecond steps, so the two require timescales that are difficult to combine in one
simulation~\parencite{aimoneComputational2016}.

A more recent theoretical treatment attributes the functional switch to the GABAergic input onto the maturing cell. While that
input is depolarizing, cooperativity drives synaptic growth so that the newborn cell becomes responsive to stimuli resembling
those that already activate mature cells; once the input becomes inhibitory, competition pushes its synaptic configuration toward
stimuli unlike those already stored. The model thereby accounts both for how adult-born cells integrate into an existing network
and for why they promote the separation of similar but not of distinct patterns~\parencite{gozelFunctional2021}, which is the same
asymmetry found experimentally in CA3~\parencite{niiboriSuppression2012}.

Reviews of this literature have noted that models agree that neurogenesis and pattern separation are related but disagree about
what is being separated, and that several distinct circuit mechanisms could produce the same behavioral
outcome~\parencite{aimoneComputational2016, wangEffect2024}. Three candidates have been set out explicitly. The first is a
population effect, in which the new cells themselves represent new information and their distinct connections to CA3 reduce
interference between memories. The second is that new cells restrain activity in the surrounding DG, making the activity of the
mature population sparser and thereby more orthogonal; the loss-of-function evidence for this is the increase in dentate and hilar
gamma bursting that follows ablation of the young population~\parencite{lacefieldEffects2012}. The third is that they
differentially engage CA3 interneurons, recruiting feedforward inhibition in CA3 and priming the recurrent
network~\parencite{aimoneComputational2016}. The first is an encoding mechanism and the other two are circuit mechanisms that act
on the mature population; distinguishing between them requires a model that includes CA3, which the studies reviewed below do not.

In a conductance-based DG model~\parencite{kimEffect2024}, the large population of mGCs performed excellent pattern separation, while
the minority of iGCs served as pattern integrators. The presence of iGCs deteriorated the DG's pattern separation capacity
compared to a DG composed solely of mGCs; however, the authors postulate that this mixed coding of pattern integration and
separation could increase the capacity for memory storage and retrieval, a hypothesis they were unable to test since only the DG
was modeled, without CA1 or CA3. Interestingly, these results contradict experimental studies demonstrating that an increase in
neurogenesis enhances pattern separation~\parencite{sahayIncreasing2011}, which may indicate that computational models of the
hippocampus have not yet fully captured the computations it performs.

In contrast, the work of~\textcite{yangDynamic2025} employed a DG model to investigate the dynamic impact of neurogenesis on
pattern separation. Their results indicate that the general effect of iGCs on GCs is inhibitory, which increases network
sparsity and consequently improves pattern separation, providing computational evidence for the ``indirect coding'' hypothesis.
The authors also investigated how iGC characteristics during maturation, such as decreasing activity and increasing synaptic
plasticity, dynamically affect pattern separation. The study further suggests that the presence of iGCs increases the
adaptability of the DG network to different input frequencies, improving the robustness of pattern separation.

The model of~\textcite{aimoneComputational2009} was implemented at an intermediate level of biological
plausibility\footnote{This model is considered to be at an intermediate level of biological plausibility because it does not
model the electrical dynamics of individual neurons through conductance-based models, as in this project, but instead uses a
digitized spike-rate model with a time step of \SI{25}{\milli\second}.}, which allowed for the investigation of iGCs throughout
their entire maturation over time. By modeling the maturation timeline of this iGC population, the model showed that iGCs
preferentially encoded information received when they were young, thus exhibiting the role of temporal integrators. This model
was corroborated by experimental studies~\parencite{berdugo-vegaSharpening2023}, but did not analyze the implications for other
areas such as CA1 and CA3, modeling only the DG\@.
